% !TEX program = xelatex
% ============================================================
%  Arabic Conference Paper Template (Short Paper) — Nibras Center
%  Compile with:
%    xelatex main.tex
%    bibtex main
%    xelatex main.tex
%    xelatex main.tex
% ============================================================
%  Features:
%    - A4, 10pt, article class, twocolumn
%    - Compact title block spanning both columns
%    - Abstract + keywords
%    - Sections: intro, related work, method, results, conclusion
%    - BibTeX bibliography
%    - Target: 2 pages
% ============================================================

\documentclass[10pt,a4paper,twocolumn]{article}

% ---- Geometry ----
\usepackage[a4paper,margin=1.8cm,top=2cm,bottom=2cm,columnsep=0.7cm,headheight=12pt]{geometry}

% ---- Core packages (order matters for bidi) ----
\usepackage{url}
\usepackage{amsmath}
\usepackage{amssymb}
\usepackage{graphicx}
\usepackage{array}
\usepackage{booktabs}
\usepackage{tabularx}
\usepackage{enumitem}
\usepackage{float}
\usepackage{caption}
\usepackage{titlesec}
\usepackage{fancyhdr}
\usepackage{setspace}

% ---- RTL / Arabic language support ----
\usepackage{bidi}
\usepackage[no-math]{fontspec}
\usepackage[quiet]{polyglossia}

% ---- Fonts ----
\setmainfont[Scale=1]{Amiri-Regular.ttf}
\newfontfamily\arabicfont[Script=Arabic,Scale=0.92]{Amiri-Regular.ttf}
\newfontfamily\englishfont[Scale=0.88]{TeX Gyre Termes}
\setmonofont{Amiri-Regular.ttf}

% ---- Languages ----
\setdefaultlanguage[calendar=gregorian,hijricorrection=1]{arabic}
\setotherlanguage[variant=british]{english}

% ---- Hyperref ----
\usepackage[breaklinks,hidelinks]{hyperref}
\hypersetup{
  pdftitle={ورقة مؤتمر},
  pdfauthor={اسم المؤلف}
}

% ---- Captions in Arabic ----
\addto\captionsarabic{%
  \renewcommand{\contentsname}{الفهرس}%
  \renewcommand{\figurename}{الشكل}%
  \renewcommand{\tablename}{الجدول}%
  \renewcommand{\refname}{المراجع}%
}

% ---- Section formatting (very compact for 2-page paper) ----
\titleformat{\section}{\normalsize\bfseries}{\thesection}{0.4em}{}
\titleformat{\subsection}{\small\bfseries}{\thesubsection}{0.3em}{}
\titlespacing*{\section}{0pt}{0.6ex plus 0.2ex}{0.4ex plus 0.1ex}
\titlespacing*{\subsection}{0pt}{0.4ex plus 0.1ex}{0.3ex}

% ---- Headers / footers ----
\pagestyle{fancy}
\fancyhf{}
\fancyhead[R]{\small\itshape عنوان مختصر}
\fancyhead[L]{\small المؤتمر الوطني ---}
\fancyfoot[C]{\small\thepage}
\renewcommand{\headrulewidth}{0.3pt}
\renewcommand{\footrulewidth}{0pt}

% ---- List spacing (very compact) ----
\setlist[itemize]{topsep=0.1em, itemsep=0.05em, parsep=0pt, leftmargin=1.2em}
\setlist[enumerate]{topsep=0.1em, itemsep=0.05em, parsep=0pt, leftmargin=1.2em}

% ---- Table column types ----
\newcolumntype{L}[1]{>{\raggedright\arraybackslash}p{#1}}
\newcolumntype{C}[1]{>{\centering\arraybackslash}p{#1}}

% ---- Line spacing (compact) ----
\setstretch{1.15}

\begin{document}

% ============================================================
% TITLE BLOCK (spans both columns)
% ============================================================
\twocolumn[
  \vspace*{0.5em}
  \begin{center}
  {\Large\bfseries عنوان ورقة المؤتمر}\par
  \vspace{0.5em}
  {\small د. اسم المؤلف$^{1}$ \quad د. اسم المؤلف الثاني$^{2}$}\par
  \vspace{0.2em}
  {\footnotesize $^{1}$القسم، الجامعة الأولى \quad $^{2}$القسم، الجامعة الثانية}\par
  \vspace{0.2em}
  {\footnotesize \texttt{email@example.com}}
  \end{center}
  \vspace{0.3em}
  \begin{center}
  \begin{minipage}{0.92\textwidth}
  \small
  \noindent\textbf{الملخص:}
  % TODO: اكتب ملخص الورقة هنا (100-200 كلمة للورقة القصيرة).
  هذا نص تجريبي للملخص. يصف المشكلة والمنهجية والنتائج الرئيسية لورقة المؤتمر القصيرة.

  \vspace{0.2em}
  \noindent\textbf{الكلمات المفتاحية:} كلمة 1، كلمة 2، كلمة 3، كلمة 4.
  \end{minipage}
  \end{center}
  \vspace{1em}
]

% ============================================================
% BODY
% ============================================================

\section{المقدمة}
\label{sec:intro}

% TODO: اكتب مقدمة موجزة (نصف عمود).
هذا نص تجريبي للمقدمة. يقدم المشكلة البحثية وأهميتها بشكل مختصر.

\section{الأعمال ذات الصلة}
\label{sec:related}

% TODO: اكتب مراجعة موجزة للأدبيات.
هذا نص تجريبي للأعمال ذات الصلة. يستعرض أبرز الأبحاث السابقة.

\section{المنهجية}
\label{sec:method}

% TODO: اكتب المنهجية بشكل مختصر.
هذا نص تجريبي للمنهجية. يصف خطوات البحث بشكل موجز.

\begin{figure}[H]
\centering
% \includegraphics[width=\columnwidth]{diagram.png}
\fbox{\parbox[c][3cm][c]{0.9\columnwidth}{\centering مكان الشكل}}
\caption{مخطط المنهجية}
\label{fig:method}
\end{figure}

\section{النتائج}
\label{sec:results}

% TODO: اكتب النتائج بشكل مختصر.
هذا نص تجريبي للنتائج. يعرض أبرز النتائج التجريبية.

\begin{table}[H]
\centering
\small
\caption{النتائج التجريبية}
\begin{tabular}{|c|c|c|}
\hline
\textbf{الطريقة} & \textbf{الدقة} & \textbf{الزمن} \\
\hline
\LR{SVM} & 85.3\% & 12ms \\
\LR{CNN} & 92.1\% & 45ms \\
\hline
\end{tabular}
\end{table}

\section{الاستنتاج}
\label{sec:conclusion}

% TODO: اكتب الاستنتاج بشكل مختصر.
هذا نص تجريبي للاستنتاج. يلخص المساهمة الرئيسية ويقترح أعمالاً مستقبلية.

% ============================================================
% BIBLIOGRAPHY
% ============================================================
\bibliography{bibliography}
\bibliographystyle{apalike}

\end{document}
